\section{Case studies}\label{sec:real}

\subsection{Data and protocol}

We use three public data sets obtained from the Rdatasets collection~\cite{rdatasets}.

\paragraph{NYC weather.} Hourly weather records at the three New York City airports (EWR, JFK, LGA) in 2013, distributed with the \texttt{nycflights13} package~\cite{nycflights13} ($n=26{,}110$ after removing five records with missing temperature, humidity or wind speed). The task is to classify whether precipitation was recorded in the hour, from temperature, dew point, humidity, wind direction and speed, pressure, month, hour and airport; $6.7\%$ of hours are positive. The loss is the logarithmic loss.

\paragraph{Melbourne pedestrians.} Hourly pedestrian counts at four sensors in Melbourne in 2015--2016, distributed with the \texttt{tsibble} package~\cite{wang2020tsibble} ($n=66{,}037$). The task is to predict $\log(1+\text{count})$ from hour, day of week, month, day of year, year and sensor, under squared loss.

\paragraph{Las Rosas maize.} Yield-monitor data from a maize field in Las Rosas, Argentina, in 1999 and 2001~\cite{anselin2004}, distributed with the \texttt{agridat} package ($n=3443$, field of about $660\times380$\,m). The task is to predict yield from the nitrogen rate, brightness value, topographic zone, year and coordinates, under squared loss.

For the two hourly data sets, blocks are two-week windows ($26$ and $53$ blocks); as the airports lie within about $35$\,km of one another and the sensors within about $2$\,km, they are treated as co-located, so that all sites observed at the same time may be correlated, and a $24$-hour buffer is removed from the training set on each side of a test block. For the maize data, blocks are $75\times75$\,m tiles ($35$ non-empty blocks) with a $10$\,m buffer. We compare XGBoost ($300$ trees of depth $6$, learning rate $0.05$, row and column subsampling $0.8$) with a random forest~\cite{breiman2001} ($200$ trees, minimum leaf size $5$, $50\%$ of features per split) and with a linear baseline (logistic or ridge regression on standardised, mean-imputed features). All models use scikit-learn~\cite{pedregosa2011} or the XGBoost library with default settings otherwise; no hyperparameter was tuned, since our object is the comparison, not the models.

\paragraph{Seed-null diagnostic.} For each data set we also compare XGBoost with itself, the two runs differing only in their random seed ($30$ pairs). By symmetry, the expected loss differential of such a comparison is zero, while in each run the stochastic subsampling produces a non-zero observed difference. A procedure that frequently declares one seed ``significantly better'' than the other confuses training noise with a real improvement. Because the data are held fixed, this diagnostic measures the sensitivity of each procedure to training randomness only; it is not a full assessment of size, which Sections~\ref{sec:sim1} and~\ref{sec:sim2} provide.

%%REALRESULTS%%
